\chapter{Systèmes post-marché}\label{ch:pl:post-trade-systems}

Le 28 mai 2024, la plupart des transactions sur titres américains ont commencé à être réglées un jour ouvré après la date de transaction au lieu de deux. La règle qui a rendu ce changement possible
exigeait aussi des courtiers qu'ils disposent d'accords ou de politiques visant à achever les affectations, les confirmations et les \omterm{def:pl:post-trade-systems:confirmation}{affirmations} avant la fin de la
date de transaction, et le
plan du secteur a avancé l'heure limite d'\omterm{def:pl:post-trade-systems:confirmation}{affirmation} du service central d'appariement de 11~h~30 le lendemain de la transaction à 21~h le jour
même de la transaction. Pour l'équipe des opérations qui corrige une transaction dont la confirmation ne concorde pas avec l'enregistrement de l'établissement, le temps disponible est passé d'une
journée ouvrée à une soirée. Ce chapitre construit les systèmes qui décident du nombre de transactions à corriger --- \omterm{def:pl:post-trade-systems:matching}{rapprochement des confirmations}, \omterm{def:pl:post-trade-systems:ssi}{instructions de règlement},
\omterm{def:pl:post-trade-systems:recon}{rapprochement} --- et mesure ce que le cycle de règlement-livraison plus court fait à celles qui restent.

\section{De la transaction enregistrée à la transaction réglée}

Le chapitre~\ref{ch:pl:trade-capture-and-booking} se terminait par une transaction enregistrée dans le modèle canonique. Entre l'enregistrement et le règlement-livraison (livre~1,
chapitre~5), la transaction doit être convenue avec la contrepartie, instruite auprès des agents qui font circuler les titres et les espèces, réglée, puis
vérifiée par rapport à ce que le dépositaire détient réellement. Chaque étape compare deux enregistrements d'une même chose, et chaque comparaison qui échoue représente
un travail manuel assorti d'une échéance.

\begin{definition}[Confirmation de transaction, affirmation]\label{def:pl:post-trade-systems:confirmation}
Une \emph{confirmation de transaction}\index{confirmation de transaction} est la déclaration, faite par la partie exécutante, des termes d'une transaction --- titre, sens, quantité, prix,
dates de transaction et de règlement-livraison, comptes --- envoyée à sa contrepartie. L'\emph{affirmation}\index{affirmation} est l'accord de la contrepartie (pour une
transaction institutionnelle, celui du gestionnaire d'actifs ou de son dépositaire) attestant que la confirmation est exacte, après quoi la transaction peut être
réglée sans autre intervention.
\end{definition}

\begin{definition}[Instruction de règlement, instruction de règlement permanente]\label{def:pl:post-trade-systems:ssi}
Une \emph{instruction de règlement}\index{instruction de règlement} donne à un agent de règlement l'ordre de livrer ou de recevoir un titre contre paiement, avec
la quantité, le montant, la date de règlement et les comptes des deux parties. Une
\emph{instruction de règlement permanente}\index{instruction de règlement permanente} (SSI) est l'ensemble des coordonnées de compte qu'une contrepartie utilise pour un marché donné, conservé une fois pour toutes comme
données de référence afin que l'instruction de chaque transaction soit générée plutôt que saisie.
\end{definition}

\begin{omfigure}
\begin{tikzpicture}[font=\small,node distance=4mm,
  box/.style={draw=omDef,fill=omDef!8,rounded corners=2pt,align=center,minimum height=10mm,text width=20mm},
  chk/.style={draw=omThm,fill=omThm!8,rounded corners=2pt,align=center,minimum height=10mm,text width=20mm},
  arr/.style={-{Stealth[length=2mm]},thick}]
\node[box] (b) {transaction enregistrée\\(ch.~21)};
\node[chk,right=of b] (m) {rappro\-chement avec la\\confir\-mation};
\node[box,right=of m] (a) {affirmation\\(heure limite)};
\node[box,right=of a] (i) {instruction\\(table des SSI)};
\node[box,right=of i] (s) {règlement\\(DVP au DCT)};
\node[chk,below=9mm of s,text width=26mm] (r) {rapprochement avec le\\dépositaire};
\node[chk,below=9mm of m,text width=34mm] (q) {file des écarts~: classer,\\attribuer, corriger, suivre l'ancienneté};
\draw[arr] (b) -- (m); \draw[arr] (m) -- (a); \draw[arr] (a) -- (i); \draw[arr] (i) -- (s); \draw[arr] (s) -- (r);
\draw[arr,omThm] (m) -- node[left,font=\scriptsize]{non-concordance} (q);
\draw[arr,omThm] (r) -- node[above,font=\scriptsize]{différence} (q);
\draw[arr,omThm,dashed] (q.east) to[bend right=12] node[below right,font=\scriptsize,pos=0.35]{corrigé} (a.south);
\end{tikzpicture}

\omcaption{Le circuit post-marché de ce chapitre~: chaque transaction enregistrée est rapprochée de la confirmation de la contrepartie, affirmée avant la
limite, instruite à partir de la table des SSI et réglée en livraison contre paiement~; les positions sont ensuite rapprochées du dépositaire. Chaque non-concordance ou
différence (en rouge) devient un écart dans une même file, où il est classé, pris en charge, corrigé et suivi en ancienneté.}\label{fig:pl:post-trade-systems:pipeline}
\end{omfigure}

La semaine du chapitre est un modèle, et chaque chiffre qu'il contient vient du modèle~: 500 transactions institutionnelles sur actions par jour, exécutées par huit courtiers pour
cinq fonds, chaque transaction étant enregistrée dans le modèle canonique de l'établissement. Les confirmations des courtiers comportent des écarts injectés à des taux donnés~: un
prix erroné pour 1~\% des transactions, une date de règlement erronée pour 1~\%, un compte erroné pour 1.5~\%, une quantité erronée pour 0.5~\%, et aucune
confirmation pour 1~\%. Trois des huit courtiers arrondissent aussi le prix moyen à quatre décimales, ce qui n'est nullement une erreur.

\section{Confirmation et affirmation}

\begin{definition}[Rapprochement des confirmations, tolérance de rapprochement]\label{def:pl:post-trade-systems:matching}
Le \emph{rapprochement des confirmations}\index{rapprochement des confirmations} associe chaque transaction enregistrée à une confirmation qui concorde sur les champs clés (titre,
sens, contrepartie, date de transaction, titulaire du compte) puis compare les autres champs~; une \emph{tolérance de rapprochement}\index{tolérance de rapprochement} est
la plus grande différence sur un champ comparé qui compte encore comme une concordance. Une paire comprise dans toutes les tolérances est rapprochée~; une paire hors de certaines d'entre elles est un
\omterm{def:pl:post-trade-systems:recon}{rapprochement} partiel, signalé avec les champs qui diffèrent~; une transaction sans candidat reste non rapprochée.
\end{definition}

\omcode{code/firm/posttrade/firm_posttrade.py}{43}{65}{Rapprochement des confirmations~: candidats par clé, meilleur candidat selon le moins de champs divergents,
et rapprochement partiel accompagné de son score et des champs qui diffèrent.}\label{lst:pl:post-trade-systems:match}

Les champs clés doivent être exacts, car ce sont eux qui décident quel enregistrement est comparé à quel autre~; les tolérances s'appliquent à ce qui est comparé. Les quantités,
les dates et les comptes n'ont aucune tolérance~: une transaction portant sur 9\,500 actions n'est pas une transaction portant sur 9\,400. Le prix est le seul champ où une tolérance
a un sens, parce que l'établissement et le courtier calculent un prix moyen sur de nombreuses exécutions et peuvent le reporter avec une précision différente.

\begin{omfigure}
\begin{tikzpicture}
\begin{semilogxaxis}[width=11.5cm,height=5.6cm,xlabel={tolérance sur le prix (\$ par action, échelle logarithmique)},ylabel={transactions rapprochées automatiquement (\%)},
  grid=major,grid style={gray!20},tick label style={font=\small},label style={font=\small},table/col sep=comma,
  ymin=55,ymax=104,ytick={60,70,80,90,100},xmin=5e-8,xmax=0.1]
\addplot[omDef,thick,mark=*] table[x=tolerance,y=rate_pct]{figdata/platforms/22-post-trade-systems/match_rate.csv};
\node[omDef,font=\scriptsize,anchor=south] at (axis cs:1e-7,60.4) {60.4};
\node[omDef,font=\scriptsize,anchor=south] at (axis cs:1e-6,61.0) {61.0};
\node[omDef,font=\scriptsize,anchor=south east] at (axis cs:1e-5,67.6) {67.6};
\node[omDef,font=\scriptsize,anchor=south] at (axis cs:1e-3,95.8) {95.8};
\node[omDef,font=\scriptsize,anchor=south] at (axis cs:5e-2,96.8) {96.8};
\draw[omThm,dashed] (axis cs:5e-5,55) -- (axis cs:5e-5,100) node[pos=0.15,right,font=\scriptsize,align=left]{une demi-unité en\\quatrième décimale};
\end{semilogxaxis}
\end{tikzpicture}

\omcaption{Part des 500 transactions du lundi rapprochées automatiquement selon la tolérance sur le prix. En dessous de 0.00005~\$, les trois courtiers qui arrondissent le
prix moyen à quatre décimales échouent sur presque toutes les transactions (60.4~\% rapprochées à 0.0000001~\$)~; de 0.00005~\$ à 0.01~\$, seuls les écarts injectés
subsistent (95.8~\%)~; à 0.05~\$, cinq prix erronés sont acceptés comme \omterm{def:pl:post-trade-systems:recon}{rapprochements} (96.8~\%). Données~:
\texttt{fig\_posttrade.py}.}\label{fig:pl:post-trade-systems:tolerance}
\end{omfigure}

La \Cref{fig:pl:post-trade-systems:tolerance} montre les deux façons de mal régler une tolérance. Trop serrée, la convention d'arrondi de trois courtiers
produit 177 écarts le lundi que personne n'a besoin de corriger~: un taux de \omterm{def:pl:post-trade-systems:recon}{rapprochement} strict de 60.4~\%. Trop lâche, un prix erroné de quelques centimes est rapproché, et
l'erreur n'est repérée --- si elle l'est --- qu'au moment où les espèces ne concordent pas. Entre les deux s'étend un plateau, de la demi-unité en quatrième décimale à un
centime, où le taux de \omterm{def:pl:post-trade-systems:recon}{rapprochement} est de 95.8~\% et où chaque écart restant est un écart réel. La bonne tolérance est la borne inférieure du plateau, dérivée de la
précision que chaque contrepartie est réputée utiliser, et consignée par contrepartie dans les données de référence plutôt que choisie globalement.

\section{Instructions de règlement}

Une fois affirmée, chaque transaction devient une instruction~: pour la première transaction du lundi, une vente de 9\,400 actions à 97.586941~\$, une \emph{livraison contre
paiement} de 917\,317.25~\$ à la date de règlement, depuis le compte de conservation du fonds vers l'agent et le compte du courtier, tirés de la table des SSI.
Rien dans l'instruction n'est saisi à la main~; chaque champ provient de la transaction canonique ou des données de référence. C'est là tout l'intérêt de la table des SSI~: un
compte erroné sur une confirmation est un écart sur une seule transaction, mais un compte erroné dans la table des SSI est un échec de règlement sur chaque transaction conclue avec cette
contrepartie sur ce marché tant que personne ne le remarque, si bien que les modifications de SSI suivent la même procédure d'approbation que les modifications de paramètres du
chapitre~\ref{ch:pl:signal-serving-and-parameter-management} --- proposée, vérifiée par une seconde personne, en vigueur à partir d'une date annoncée, jamais
modifiée sur place.

Dans le modèle, les mauvais comptes sont l'écart le plus fréquent~: 43 des 138 écarts de la semaine, contre 32 confirmations manquantes, 27
dates de règlement erronées, 26 prix erronés et 10 quantités erronées (\cref{fig:pl:post-trade-systems:kinds}). C'est aussi le type d'écart qu'une société peut
supprimer à la source, en tenant les SSI à jour une seule fois et en les partageant avec ses contreparties, au lieu de laisser chaque partie les saisir.

\section{Rapprochement et écarts}

\begin{definition}[Rapprochement, écart de rapprochement]\label{def:pl:post-trade-systems:recon}
Le \emph{rapprochement}\index{rapprochement} est la comparaison périodique de deux relevés qui doivent concorder~--- les positions ou la trésorerie de la société
face au relevé d'un dépositaire ou d'un prime broker, les transactions face à celles du courtier compensateur~--- élément par élément. Un \emph{écart
de rapprochement}\index{ecart@écart de rapprochement} est un élément pour lequel les deux relevés divergent~: présent d'un seul côté, ou présent des deux côtés avec des valeurs différentes.
\end{definition}

\omcode{code/firm/posttrade/firm_posttrade.py}{106}{120}{Rapprochement par clé, avec des groupes un-à-plusieurs~: la position globale d'un dépositaire est
comparée à la somme des allocations de la société qui s'y rattachent.}\label{lst:pl:post-trade-systems:recon}

La difficulté du \omterm{def:pl:post-trade-systems:recon}{rapprochement} tient rarement à la comparaison~; elle tient aux clés. Le dépositaire peut déclarer une seule position là où la société en détient cinq
allocations (un groupe un-à-plusieurs), un mouvement de trésorerie peut être la somme des transactions d'une journée (un agrégat), et les identifiants peuvent différer entre les deux
côtés, si bien que les données de référence du chapitre~\ref{ch:pl:reference-data-and-symbology-services} permettent de dire si un écart est réel. Un écart a alors besoin de trois éléments de plus
pour être utile~: un type (élément manquant d'un côté, écart de quantité), un responsable et un âge.

\begin{definition}[Vieillissement des écarts]\label{def:pl:post-trade-systems:ageing}
Le \emph{vieillissement des écarts}\index{vieillissement des écarts} est la classification des écarts ouverts selon le temps écoulé depuis leur ouverture, par tranches (moins d'un jour, un
à trois jours, trois à sept jours, plus de sept jours), publiée chaque jour avec le responsable de chaque écart~; les tranches anciennes sont l'endroit où le risque se cache, car un écart
qui survit à plusieurs \omterm{def:pl:post-trade-systems:recon}{rapprochements} est un écart que personne ne comprend.
\end{definition}

\section{Fonctionner sous un cycle de règlement-livraison court}

L'équipe des opérations du modèle compte trois personnes. Les écarts sont traités par valeur décroissante, chacun prenant un temps aléatoire dont la moyenne dépend du type
(30~minutes pour une date de règlement, 45 pour un prix ou une quantité, une heure pour un compte, une heure et demie pour relancer une confirmation manquante). La
fenêtre est l'endroit où les deux cycles diffèrent. Sous T+2, le modèle accorde à l'équipe le jour ouvré suivant, de 08~h à 17~h, avant qu'une transaction ne puisse plus être
réglée à temps~; sous T+1, il lui accorde la soirée de la date de transaction, de 17~h jusqu'à la limite d'\omterm{def:pl:post-trade-systems:confirmation}{affirmation} de 21~h. Ces fenêtres sont des choix du modèle,
fixés à partir des heures limites citées dans l'accroche~; une équipe réelle travaille aussi pendant la séance de négociation, ce qui profite également aux deux cycles.

\omcode{code/platforms/22-post-trade-systems/python/pl_posttrade.py}{93}{109}{L'équipe traite les écarts par valeur décroissante~; un écart
qui ne peut être clos dans la fenêtre n'est pas réglé.}\label{lst:pl:post-trade-systems:resolve}

\begin{table}[htbp]
\centering\small
\begin{tabular}{@{}>{\raggedright\arraybackslash}p{1.2cm}>{\centering\arraybackslash}p{1.5cm}>{\centering\arraybackslash}p{1.5cm}>{\centering\arraybackslash}p{1.5cm}>{\centering\arraybackslash}p{1.5cm}>{\centering\arraybackslash}p{1.5cm}>{\centering\arraybackslash}p{1.5cm}>{\centering\arraybackslash}p{1.5cm}@{}}
\toprule
cycle & écarts & résolus & en suspens & valeur en suspens & \omterm{def:pl:post-trade-systems:recon}{écarts de rapprochement} & âgés de 1 à 3 jours & âgés de 3 à 7 jours\\
\midrule
T+2 & 138 & 128 & 10 & 2.74~M\$ & 4 & 4 & 0\\
T+1 & 138 & 91 & 47 & 19.17~M\$ & 24 & 12 & 5\\
\bottomrule
\end{tabular}
\caption{La même semaine, soit 2\,500 transactions et 138 écarts sous les deux cycles, avec trois personnes aux opérations. Le \omterm{def:pl:post-trade-systems:recon}{rapprochement} avec le dépositaire
le vendredi soir porte sur les transactions à régler d'ici le vendredi~; sous T+1, 7 des 24 écarts ont moins d'un jour.}
\label{tab:pl:post-trade-systems:cycles}
\end{table}

Les écarts sont les mêmes sous les deux cycles~--- ils viennent des confirmations, et non du calendrier~--- mais le travail qu'ils demandent ne tient pas dans la
soirée. Les écarts de la semaine demandent en moyenne environ 26 heures de travail par jour~; le jour ouvré suivant donne à trois personnes 27 heures, la soirée 12. Sous
T+2, dix transactions sont en suspens sur la semaine, pour 2.74~millions de dollars de valeur~; sous T+1, quarante-sept le sont, pour 19.17~millions de dollars, et le \omterm{def:pl:post-trade-systems:recon}{rapprochement} du vendredi trouve six
fois plus d'écarts de positions (\cref{tab:pl:post-trade-systems:cycles}). Une transaction en suspens n'est pas perdue~--- elle se règle plus tard~--- mais le fonds
ne dispose pas des titres ni de la trésorerie qu'il attendait~; le \omterm{def:pl:position-and-pnl-services:service}{service de positions} du chapitre~\ref{ch:pl:position-and-pnl-services} doit tenir séparément les positions réglées et
non réglées, et certains marchés font payer le suspens~: le livre~2, chapitre~5, décrit la pénalité de suspens sur le marché du Trésor.

\begin{omfigure}
\begin{tikzpicture}
\begin{axis}[width=11.5cm,height=5.4cm,ybar,bar width=9pt,ylabel={transactions},symbolic x coords={account,missing,settle_date,price,qty},
  xticklabels={mauvais\\compte,confirmation\\manquante,mauvaise\\date de\\règlement,mauvais\\prix,mauvaise\\quantité},xtick=data,
  xticklabel style={align=center,font=\footnotesize},
  grid=major,grid style={gray!20},tick label style={font=\small},label style={font=\small},table/col sep=comma,ymin=0,ymax=62,
  nodes near coords,every node near coord/.append style={font=\scriptsize},
  legend style={font=\small,at={(0.98,0.97)},anchor=north east,legend columns=1,cells={anchor=west},
    fill=white,draw=gray!40}]
\addplot[fill=omDef!35,draw=omDef] table[x=kind,y=breaks]{figdata/platforms/22-post-trade-systems/kinds.csv};
\addplot[fill=omThm!25,draw=omThm] table[x=kind,y=failed_t2]{figdata/platforms/22-post-trade-systems/kinds.csv};
\addplot[fill=omThm!70,draw=omThm] table[x=kind,y=failed_t1]{figdata/platforms/22-post-trade-systems/kinds.csv};
\legend{écarts de la semaine,en suspens sous T+2,en suspens sous T+1}
\end{axis}
\end{tikzpicture}

\omcaption{Les 138 écarts de la semaine par type, et les transactions non réglées sous chaque cycle. Les mauvais comptes et les confirmations manquantes sont
les plus fréquents et les plus lents à corriger, et ensemble ils représentent 36 des 47 suspens sous T+1. Données~: \texttt{fig\_posttrade.py}.}
\label{fig:pl:post-trade-systems:kinds}
\end{omfigure}

Deux moyens permettent de ramener les suspens sous T+1 au niveau de ceux de T+2, et la \cref{fig:pl:post-trade-systems:staffing} les compare. Le premier tient aux effectifs~:
avec sept personnes aux opérations au lieu de trois, la soirée absorbe le travail, et les suspens de la semaine tombent à 11. Le second levier consiste à réduire le nombre d'écarts~: diviser par deux
le taux de mauvais comptes et de confirmations manquantes~--- des SSI tenues à jour une fois et partagées, des allocations envoyées au courtier dès qu'elles sont
arrêtées~--- ramène les suspens à 13 avec les mêmes trois personnes. Cette seconde voie est moins coûteuse et elle tient dans la durée, ce qui explique que le passage à T+1 ait été préparé
autant dans les données de référence et les flux de travail d'allocation que dans les effectifs.

\begin{omfigure}
\begin{tikzpicture}
\begin{axis}[width=11.5cm,height=5.6cm,xlabel={personnes aux opérations},ylabel={transactions en suspens dans la semaine (T+1)},
  grid=major,grid style={gray!20},tick label style={font=\small},label style={font=\small},table/col sep=comma,
  xmin=2.6,xmax=8.6,ymin=0,ymax=52,xtick={3,4,5,6,7,8}]
\addplot[omThm,thick,mark=*] table[x=staff,y=failed]{figdata/platforms/22-post-trade-systems/staffing.csv};
\addplot[omDef,thick,mark=square*] table[x=staff,y=failed_half_errors]{figdata/platforms/22-post-trade-systems/staffing.csv};
\draw[gray,dashed] (axis cs:2.6,10) -- (axis cs:8.6,10) node[pos=0.97,below left,font=\scriptsize]{T+2, trois personnes~: 10};
\node[font=\scriptsize,anchor=north east,align=right] at (axis cs:8.5,50) {\textcolor{omThm}{$\bullet$ taux d'erreur du modèle}\\
  \textcolor{omDef}{$\blacksquare$ mauvais comptes et confirmations manquantes réduits de moitié}};
\end{axis}
\end{tikzpicture}

\omcaption{Transactions en suspens dans la semaine sous T+1 selon la taille de l'équipe des opérations, aux taux d'erreur du modèle (en rouge) et avec les
mauvais comptes et confirmations manquantes réduits de moitié (en bleu). Sept personnes ramènent les suspens à 11~; réduire de moitié les deux erreurs les plus fréquentes les ramène à 13
avec trois personnes. Données~: \texttt{fig\_posttrade.py}.}\label{fig:pl:post-trade-systems:staffing}
\end{omfigure}

La leçon générale est celle du chapitre~\ref{ch:pl:trade-capture-and-booking}~: le coût du post-trade se fixe en amont. Chaque écart que le pipeline
crée est payé par une personne sous la contrainte d'une échéance, et un cycle plus court raccourcit l'échéance sans changer le nombre d'écarts.
La mesure à surveiller est la part des transactions qui vont de l'enregistrement au règlement-livraison sans intervention manuelle~--- le taux de \omterm{def:pl:trade-capture-and-booking:stp}{traitement de bout en bout}~--- par
contrepartie et par type d'écart, ainsi que l'âge de ce qui reste.

\begin{dated}{2026-09}{Le passage des États-Unis à T+1}\label{dat:pl:post-trade-systems:t1}
La SEC Release 34-96930 (2023) a modifié la Rule 15c6-1 pour raccourcir le cycle de règlement-livraison standard de la plupart des transactions effectuées par les broker-dealers, de T+2 à T+1,
avec une date d'entrée en application du 28 mai 2024, le lendemain d'un jour férié fédéral. La nouvelle Rule 15c6-2 oblige les courtiers à disposer d'accords ou de politiques visant
à finaliser les allocations, les confirmations et les \omterm{def:pl:post-trade-systems:confirmation}{affirmations} dès que cela est techniquement possible, et au plus tard à la fin de la date de transaction~;
la nouvelle Rule 17Ad-27 oblige les chambres de compensation qui fournissent un service central de \omterm{def:pl:post-trade-systems:matching}{rapprochement des confirmations} à faciliter le \omterm{def:pl:trade-capture-and-booking:stp}{traitement de bout en bout} et à en rendre compte
chaque année. Le texte cite le T+1 Report du secteur, qui a avancé la limite d'\omterm{def:pl:post-trade-systems:confirmation}{affirmation} de 11~h~30, heure de l'Est, le lendemain de la transaction
à 21~h, heure de l'Est, le jour même de la transaction.
\end{dated}

\section{Tutoriel~: une semaine de confirmations sous deux cycles}\label{tut:pl:post-trade-systems:tut}

\begin{tutorial}
\textbf{Objectif.} Rapprocher une semaine de transactions comportant des écarts introduits volontairement, donner les instructions, réconcilier et dater les écarts, puis rejouer la semaine en T+1.
\textbf{État final~:} la \cref{fig:pl:post-trade-systems:tolerance} et le \cref{tab:pl:post-trade-systems:cycles}.
\begin{tutsteps}
\item \textbf{Transactions et confirmations}~: \texttt{pl\_posttrade.trades(0, 1)} et \texttt{confirms(ours)}~; compter les types d'écarts introduits.
\item \textbf{Appariement}~: \texttt{match\_rates(tols=\dots)} pour des tolérances allant de $10^{-7}$ à $0.05$.
\item \textbf{Instructions}~: construire une \texttt{SSITable} pour les huit courtiers et générer l'instruction de chaque transaction appariée.
\item \textbf{Écarts et corrections}~: \texttt{week(\textquotedbl T+2\textquotedbl)} et \texttt{week(\textquotedbl T+1\textquotedbl)}.
\item \textbf{Réconciliation}~: \texttt{custodian\_recon(cycle)}, les écarts et leur vieillissement.
\end{tutsteps}
\textbf{Pour aller plus loin.} Ajouter une réconciliation de trésorerie (les paiements nets du jour par fonds, comparés au relevé de trésorerie du dépositaire) avec
un appariement agrégé~; laisser l'équipe travailler une partie de la journée de négociation et vérifier si le T+1 exige encore davantage de personnel.
\end{tutorial}

\section{Construction~: post-marché}\label{bld:pl:post-trade-systems:posttrade}

\begin{build}
\bfield{Rôle} Faire accepter chaque transaction par sa contrepartie, la donner en instruction et démontrer ensuite que les livres concordent avec le dépositaire, avec chaque
exception dans une seule file d'attente.
\bfield{Interface} \texttt{match}, \texttt{match\_rate}, \texttt{MatchResult}, \texttt{SSITable.instruction}, \texttt{reconcile} (avec
\texttt{groups}), \texttt{Break}, \texttt{classify}, \texttt{age}.
\bfield{Règles} Clés exactes, tolérances par champ et consignées~; instructions générées à partir des SSI, jamais saisies à la main~; modifications de SSI approuvées et datées~;
chaque écart a un type, un responsable et une heure d'ouverture~; vieillissement communiqué chaque jour.
\bfield{Tests d'acceptation} \texttt{code/firm/posttrade/tests/}~: transactions appariées, partielles et non appariées des deux côtés~; le meilleur candidat choisi
parmi plusieurs~; une instruction issue de la table des SSI et une SSI manquante~; réconciliation un-à-plusieurs et les tranches de vieillissement.
\bfield{Pour aller plus loin} Réconciliation de trésorerie et de transactions avec appariement agrégé~; un flux d'\omterm{def:pl:post-trade-systems:confirmation}{affirmation} avec heures limites et escalade~; un
rapport de taux de \omterm{def:pl:trade-capture-and-booking:stp}{traitement de bout en bout} par contrepartie.
\end{build}

\begin{omsources}
\item US Securities and Exchange Commission, \emph{Shortening the Securities Transaction Settlement Cycle}, Release No.~34-96930, 2023.
\item One Quant Book~1, chapitre~5 (règlement-livraison, cycle de règlement-livraison, suspens, livraison contre paiement)~; livre~2, chapitre~5 (la pénalité de suspens)~;
livre~16, chapitre~4 (le dépositaire).
\end{omsources}

\section{Exercices}

\begin{exercise}[$\star$]\label{exo:pl:post-trade-systems:1}
Pourquoi les champs clés sont-ils appariés exactement et seuls les autres avec une tolérance~?
\end{exercise}

\begin{exercise}[$\star$]\label{exo:pl:post-trade-systems:2}
Un courtier arrondit les prix à quatre décimales. Quelle est la plus petite tolérance de prix pour laquelle ses confirmations correctes s'apparient toujours~?
\end{exercise}

\begin{exercise}[$\star$]\label{exo:pl:post-trade-systems:3}
En T+1, combien d'écarts de la semaine ont été corrigés, et quelle part des écarts a donné lieu à un suspens~?
\end{exercise}

\begin{exercise}[$\star\star$]\label{exo:pl:post-trade-systems:4}
Pourquoi un mauvais compte dans la table des SSI est-il plus grave qu'un mauvais compte sur une seule confirmation~?
\end{exercise}

\begin{exercise}[$\star\star$]\label{exo:pl:post-trade-systems:5}
Le dépositaire signale une position de 3\,000 actions alors que la société détient des allocations de 1\,000 et de 2\,000 à deux fonds sous le même compte.
Comment la réconciliation évite-t-elle trois écarts~?
\end{exercise}

\begin{exercise}[$\star\star$]\label{exo:pl:post-trade-systems:6}
Pourquoi la réconciliation du vendredi ne trouve-t-elle que 4 écarts de position en T+2 alors que 10 transactions ont fait l'objet d'un suspens~?
\end{exercise}

\begin{exercise}[$\star\star\star$]\label{exo:pl:post-trade-systems:7}
\emph{Programmation.} Relancer \texttt{week(\textquotedbl T+1\textquotedbl, staff=s)} pour $s$ allant de 3 à 8 et trouver la plus petite équipe qui ne laisse pas en suspens plus de
transactions qu'en T+2 avec trois personnes.
\end{exercise}

\begin{exercise}[$\star\star\star$]\label{exo:pl:post-trade-systems:8}
\emph{Trouver la faille.} «~Notre taux d'appariement est passé de 95.8~\% à 96.8~\% ce mois-ci depuis que nous avons élargi la tolérance de prix à cinq centimes~:
les opérations font un meilleur travail.~»
\end{exercise}

\section{Problème~: une soirée au lieu d'une journée}

\begin{problem}[{Problème du week-end --- une soirée au lieu de deux jours}]\label{pb:pl:post-trade-systems:1}
La semaine du chapitre~: 500 transactions par jour, huit courtiers, cinq fonds, écarts introduits volontairement, trois personnes aux opérations.

\medskip\noindent\textbf{Partie I --- Faire accepter la transaction.}
\begin{enumerate}
\item Que compare-t-on entre une transaction enregistrée et une confirmation, et que doit-on exiger d'exact~?
\item Qu'apporte l'\omterm{def:pl:post-trade-systems:confirmation}{affirmation} à l'appariement~?
\item Pourquoi trois courtiers produisent-ils des écarts avec une tolérance serrée, et combien de transactions s'apparient à 0.0000001~\$~?
\item Où se situe le palier du taux d'appariement, et quel est ce taux à cet endroit~?
\item Que se passe-t-il à une tolérance de cinq centimes~?
\end{enumerate}

\medskip\noindent\textbf{Partie II --- Donner les instructions et réconcilier.}
\begin{enumerate}[resume]
\item Que contient l'\omterm{def:pl:post-trade-systems:ssi}{instruction de règlement} de la première transaction du lundi, et d'où vient chaque champ~?
\item Comment doit-on contrôler les modifications de SSI~?
\item Qu'est-ce qui rend la réconciliation difficile en pratique~?
\item De quoi un écart a-t-il besoin en plus de la différence elle-même~?
\item Pourquoi les anciennes tranches de vieillissement sont-elles celles qu'il faut surveiller~?
\end{enumerate}

\medskip\noindent\textbf{Partie III --- Les deux cycles.}
\begin{enumerate}[resume]
\item Quelle fenêtre le modèle laisse-t-il à l'équipe dans chaque cycle, et pourquoi~?
\item Quelle charge de travail les écarts représentent-ils par jour, et combien chaque fenêtre en permet-elle~?
\item Combien de transactions font l'objet d'un suspens dans chaque cycle, et pour quelle valeur~?
\item Quels types d'écarts font le plus de suspens en T+1~?
\item Que trouve la réconciliation du vendredi dans chaque cycle~?
\end{enumerate}

\medskip\noindent\textbf{Partie IV --- Le verdict.}
\begin{enumerate}[resume]
\item Énoncer le \emph{résultat nommé}~: la part des transactions appariées automatiquement à chaque tolérance, les écarts subsistant à l'heure limite d'\omterm{def:pl:post-trade-systems:confirmation}{affirmation}
dans un cycle de deux jours et dans un cycle d'un jour, et les suspens qui en découlent.
\item Combien de personnes le T+1 exige-t-il pour ne pas laisser en suspens plus de transactions qu'en T+2 avec trois~?
\item Qu'apporte la réduction de moitié des mauvais comptes et des confirmations manquantes~?
\item Lequel des deux financeriez-vous, et pourquoi~?
\item En une phrase~: que change un cycle de règlement plus court~?
\end{enumerate}
\end{problem}

\section{Questions d'entretien}

\begin{interviewq}[$\star$ \iqroles{développeur}]\label{iq:pl:post-trade-systems:1}
Qu'est-ce qu'un \omterm{def:pl:post-trade-systems:recon}{écart de rapprochement}, et que mettriez-vous dans un rapport d'écarts~?
\end{interviewq}

\begin{interviewq}[$\star\star$ \iqroles{développeur}]\label{iq:pl:post-trade-systems:2}
Comment choisissez-vous les tolérances d'un moteur d'appariement des confirmations~?
\end{interviewq}

\begin{interviewq}[$\star\star$ \iqroles{développeur}]\label{iq:pl:post-trade-systems:3}
Votre dépositaire signale les positions par compte, et vos livres les tiennent par fonds et par stratégie. Concevez la réconciliation.
\end{interviewq}

\begin{interviewq}[$\star\star$ \iqroles{développeur}]\label{iq:pl:post-trade-systems:4}
Le marché passe à un cycle de règlement d'un jour. Que change-t-il dans vos systèmes de post-marché et dans votre équipe~?
\end{interviewq}

\begin{interviewq}[$\star\star\star$ \iqroles{développeur}]\label{iq:pl:post-trade-systems:5}
Concevez la plateforme de post-marché d'une société qui négocie des actions avec trente courtiers pour cent fonds.
\end{interviewq}
